\section{Conclusion}
\label{sec:conclusion}

The capability lattice provides the first tractable instantiation of
goal-frontier maximization. By constructing the capability space as a finite
weighted poset discovered through dual-channel discovery, the lattice resolves
the structure-of-$\G$ problem that the foundational paper identified as its
most urgent open question. The lattice measure $\volL$, defined via Möbius
inversion, satisfies the axioms that unlock all of the foundational paper's
structural results: self-balancing, scorpion detection, and cooperative
expansion hold on finite lattices exactly as they do on abstract capability
spaces.

The construction also advances the estimator. On the lattice, agent reports
carry real-valued weighted changes rather than ternary signs, resolving the
foundational paper's magnitude-estimation limitation and strengthening the
sign-correctness conditions for mixed-sign capability changes. The bootstrap
procedure shows how a GFM actor can discover the lattice incrementally through
communication and observation, with cross-validation between the two channels
providing a structured trust signal.

The lattice is bounded by what can be communicated and observed. Every
discovered capability carries at least binary weight (1 bit); graded benchmarks
refine this upward. Capabilities the actor has not encountered remain outside
the lattice, but no known capability is invisible. The gap between lattice
volume and experiential optionality narrows under the binary default but does
not close: capabilities with latent gradations are understated rather than
ignored. Developing richer benchmarks for experiential and relational
capabilities remains the central challenge for the research program.
